\documentclass[11pt,a4paper]{article}
\usepackage[T2A]{fontenc}
\usepackage[utf8]{inputenc}
\usepackage[russian,english]{babel}
\usepackage{amsmath,amssymb,amsthm,mathtools}
\usepackage{setspace}
\usepackage{enumitem}
\usepackage[hidelinks]{hyperref}
\setstretch{1.3}
\hyphenpenalty=10000
\exhyphenpenalty=10000
\renewcommand{\leq}{\leqslant}
\renewcommand{\le}{\leqslant}
\renewcommand{\geq}{\geqslant}
\renewcommand{\ge}{\geqslant}
\newcommand{\head}[1]{\par\bigskip\noindent\textbf{\underline{#1}}\quad}
\begin{document}
\begin{center}
{\huge\bfseries Integrals -- 1\par}\vspace{1.2em}
{\Large Substitutions: rationalisation, Euler, the half-angle substitution\par}\vspace{0.8em}
{\foreignlanguage{russian}{1 сентября 2026 года}\par}
\end{center}

\noindent\textbf{\underline{Theoretical minimum.}}\par\vspace{-\bigskipamount}\medskip

\head{Def.} A substitution reduces an integral to a class we can already handle. Three classes must
be recognised instantly: $R(x,x^{p_1/q_1},\dots)$ is rationalised by $x=t^N$ with $N$ the common
denominator of the exponents; $R\bigl(x,\sqrt{ax^2+bx+c}\bigr)$ is handled by Euler's
substitutions; $R(\sin x,\cos x)$ by the half-angle substitution $t=\tan\frac x2$.

\head{Theorem 1 (Euler's substitutions).} Let $y=\sqrt{ax^2+bx+c}$. If $a>0$, put
$y=\pm\sqrt a\,x+t$; if $c>0$, put $y=xt\pm\sqrt c$; if the trinomial has real roots $x_1,x_2$, put
$y=t(x-x_1)$. In each case both $x$ and $y$ become rational functions of $t$.

\head{Theorem 2 (half-angle substitution).} For $t=\tan\frac x2$ one has $\sin x=\frac{2t}{1+t^2}$,
$\cos x=\frac{1-t^2}{1+t^2}$ and $dx=\frac{2\,dt}{1+t^2}$. If the integrand is even in the pair
$(\sin x,\cos x)$, the substitution $t=\tan x$ is cheaper; if it is odd in $\sin x$, use $t=\cos x$,
and symmetrically.

\head{Fact 1.} Integration by parts is not a trick but the identity $\int u\,dv=uv-\int v\,du$: the
only question is what to call $dv$, and the answer is almost always the most unpleasant factor that
can still be integrated.

\noindent\rule{\textwidth}{0.4pt}\par\medskip

\begin{enumerate}[itemsep=0.6em]
\item Derive all three Euler substitutions and explain when each of them is the cheapest. Compute
$\int\frac{dx}{x\sqrt{x^2+4x+3}}$ and $\int\sqrt{x^2+1}\,dx$, the second one in two ways (an Euler
substitution and a hyperbolic one), and compare the answers.

\item Derive the half-angle formulas and compute $\int\frac{dx}{2+\cos x}$, $\int\frac{dx}{1+\sin x}$
(Integration Bee 2018) and $\int\frac{dx}{\sin x+\tan x}$. Explain why in the third integral
$t=\tan\frac x2$ beats the substitution $t=\cos x$.

\item (Integration Bee) Rationalising substitutions:
\begin{enumerate}[label=\alph*),itemsep=0pt,topsep=0.3em]
\item $\displaystyle\int\bigl(1+x^{1/2}+x^{1/3}\bigr)\bigl(1+x^{-1/2}+x^{-1/3}\bigr)dx$ (2016);
\item $\displaystyle\int\frac{dx}{5+4\sqrt x+x}$ (2015);
\item $\displaystyle\int\frac{x+\sqrt x}{1+\sqrt x}\,dx$ (2025);
\item $\displaystyle\int\frac{dx}{x+3\sqrt x}$ (2019).
\end{enumerate}

\item (Integration Bee, speed drill) Compute:
\begin{enumerate}[label=\alph*),itemsep=0pt,topsep=0.3em]
\item $\int x^2\sqrt{x^3+2}\,dx$ (2017);
\item $\int x^3\sqrt{x^2+1}\,dx$ (2016);
\item $\int\frac{x\,dx}{x^4+x^2+1}$ (2016);
\item $\int\frac{dx}{x(x^2+1)}$ (2017);
\item $\int\frac{x^5}{2+x^{12}}\,dx$ (2018);
\item $\int\frac{dx}{\sqrt{x\sqrt x-x^2}}$ (2018).
\end{enumerate}

\item (Integration Bee) Self-similar integrands: $\displaystyle\int\sqrt{x+\sqrt{x+\sqrt{x+\cdots}}}\,dx$
(2014), $\displaystyle\int\sqrt{x\sqrt{x\sqrt{x\cdots}}}\,dx$ (2020) and
$\displaystyle\int_0^1\sqrt{x^2+x+\sqrt{x^2+x+\sqrt{\cdots}}}\,dx$ (2026). First say precisely what
the infinite radical means, and only then integrate.
\end{enumerate}
\end{document}
